تصدرت التحديات المرتبطة \textbf{بارتفاع تكاليف الطاقة والمياه} قائمة المعوقات التي واجهت الشركات خلال الربع محل الدراسة في ظل الارتفاع المستمر في أسعارها الأمر الذي ينعكس في زيادة تكاليف الإنتاج خاصة للأنشطة كثيفة الاستهلاك للطاقة والمياه، مما يمثل عبئا إضافيا على الشركات مقارنة بالربع السابق. واحتلت التحديات المرتبطة \textbf{بإجراءات التعامل مع الجهات الحكومية} المرتبة الثانية بالرغم من التراجع الملحوظ؛ حيث تتسبب في معاناة مجتمع الأعمال نتيجة بطء الإجراءات، والروتين، وتعامل بعض الموظفين، بالإضافة إلى تعدد موظفي الضبطية القضائية في معظم الجهات الحكومية، مما يفتح مجال للفساد والمصروفات غير الرسمية. وجاءت التحديات المرتبطة بارتفاع معدلات \textbf{التضخم} في المرتبة الثالثة.
